\chapter{ಪ್ರೋಗ್ರಾಂ ಶೀರ್ಷಿಕೆ}
ಮೊದಲೇ ನೋಡಿರುವ ಅಧ್ಯಾಯಗಳಲ್ಲಿ,  ಸ್ಪ್ರಯ್ಟ್ ಅನ್ನು ಒಡಾಡಿಸಲು ಪ್ರೋಗ್ರಾಂ ಬರೆಯುವುದರಲ್ಲಿ ತಪ್ಪುಗಳಾಗುವುದು ಸಹಜ.  ಒಮ್ಮೊಮ್ಮೆ ಪ್ರೋಗ್ರಾಂ ತಪ್ಪಾದಾಗ,  ಸ್ಪ್ರಯ್ಟ್ ವೇದಿಕೆಯಿಂದ ಮಾಯವಾಗಿಯೂ ಬಿಡಬಹುದು.  ನಿಜವಾಗಿ ನೋಡಿದರೆ ಸ್ಪ್ರಯ್ಟ್ ವೇದಿಕೆಯಿಂದ ಹೊರಬಿದ್ದಿರ ಬಹುದು. ಹಾಗಾದಾಗ ಅದು ನಮಗೆ ಮಾಯವಾದಂತೆ ಕಾಣಬಹುದು.  ನಾವು ಚುರುಕಾಗಿ ನೋಡಿಕೊಂಡರೆ, ತಪ್ಪಗಿ ನೀಡಿದ ಆಜ್ಞೆಗಳೆಲ್ಲವನ್ನೂ ಕ್ರಮಾನುಸಾರವಾಗಿ ಹಿಂತೆಗೆದರೆ, ಸ್ಪ್ರಯ್ಟ್ ಪುನಃ ವೇದಿಕೆಯ ಮೇಲೆ ಕಾಣಿಸಿಕೊಳ್ಳುವುದು. ಆದರೆ ಈ ರೀತಿ ಮಾಡುವುದು ಸುಲಭವಲ್ಲ. 

ಹಾಗಾದರೆ, ಕಣ್ಮರೆಯಾದ ಸ್ಪ್ರಯ್ಟ್ ಅನ್ನು ಪುನಃ ವೇದಿಕೆಯ ಮೇಲೆ ತರಿಸಿ ಹೊಸದಾಗಿ ಪ್ರೋಗ್ರಾಂ ಪ್ರಾರಂಭಿಸುವುದು ಹೇಗೆ? ಮತ್ತೊಂದು ರೀತಿಯಲ್ಲಿ ಕೇಳುವುದಾದರೆ, ಪ್ರತಿ ಸಲವೂ ಸ್ವಚ್ಛ ವೇದಿಕೆಯಿಂದ ಪ್ರೋಗ್ರಾಂ ಪ್ರಾರಂಭಿಸುವಂತೆ ಮಾಡುವುದು ಹೇಗೆ? ಇದಕ್ಕೆ, ಸ್ಕ್ರಾಚ್ ಮಾತ್ರವಲ್ಲ, ಬೇರೆ ಎಲ್ಲಾ ಪ್ರೋಗ್ರಾಂ ಭಾಷೆಗಳಲ್ಲಿಯೂ ಪ್ರೋಗ್ರಾಂ ಶೀರ್ಷಿಕೆಗಳನ್ನು ಬಳಸುತ್ತಾರೆ.  ಸ್ಕ್ರಾಚ್ ಪ್ರೋಗ್ರಾಂ ಶೀರ್ಷಿಕೆಗಳನ್ನು ಈಗ ಉದಾಹರಣೆಗಳೊಂದಿಗೆ ಕಲಿಯೋಣ. 


\section{ಸ್ವಚ್ಛ ವೇದಿಕೆ}
\begin{figure}[h]
\begin{multicols}{3}
\begin{center}
\begin{Scratch}[1]
\scbox{ಅಳಿಸು}{pen}
\end{Scratch}

\begin{Scratch}[1]
\scbox{ಅಳಿಸು}{pen}
\scbox{\textenglish{x:} \cb[w]{0} \textenglish{y:}\cb[w]{0}  ಗೆ ಹೋಗು}{motion}
\end{Scratch}

\begin{Scratch}[1]
\scbox{ಅಳಿಸು}{pen}
\scbox{\textenglish{x:} \cb[w]{0} \textenglish{y:}\cb[w]{0} ಗೆ ಹೋಗು}{motion}
\scbox{\cb[w]{90} ಯ ದಿಕ್ಕಿನಲ್ಲಿಡು}{motion}
\end{Scratch}

\end{center}
\end{multicols}
\caption{ಸ್ಕ್ರಾಚ್ ಪ್ರೋಗ್ರಾಂ ಶೀರ್ಷಿಕೆ}
\label{hdr_init1}
\end{figure}

ಚಿತ್ರ \ref{pen_sqr}ನಲ್ಲಿನ ಎಡಭಾಗದಲ್ಲಿರುವಂತೆ ಪ್ರೋಗ್ರಾಮ್ಮನ್ನು ಮತ್ತೆ ಮೂಡಿಸಿ ಪ್ರೋಗ್ರಾಂ ಅನ್ನು  ಚಲಾಯಿಸಿ (ವೇದಿಕೆಯ ಮೇಲಿನ ಹಸಿರು ಧ್ವಜವನ್ನು ಒತ್ತಿ) ಚೌಕ ಮೂಡಿಬರುವುದನ್ನು ಗಮನಿಸಿ.  ಈಗ ಪ್ರೋಗ್ರಾಂ ಅನ್ನು  ಮತ್ತೊಮ್ಮೆ ಚಲಾಯಿಸಿ ಏನಾಗುತ್ತದೆಂದು ಗಮನಿಸಿ. ಹೀಗೆ ಎರಡು ಚೌಕಗಳೇಕೆ ಮೂಡಿಬಂದವು? ನಾವು ಪ್ರೋಗ್ರಾಂ  ಬರೆದದ್ದು ಒಂದು ಚೌಕ ಬಿಡಿಸುವುದಕಲ್ಲವೇ? ಹೀಗಾಗಬಾರದೆಂದರೆ, ನಾವು ಪ್ರೋಗ್ರಾಂ ಅನ್ನು  ಚಲಾಯಿಸುವ ಮುನ್ನ ವೇದಿಕೆಯನ್ನು ಸ್ವಚ್ಛಗೊಳಿಸಬೇಕು. ಅದಕ್ಕೆ ಸ್ಕ್ರಾಚ್ನ "ಲೇಖನಿ" ಗುಂಪಿನಲ್ಲಿ ಅಳಿಸು ಎಂಬ ಆಜ್ಞೆಯನ್ನು ಉಪಯೋಗಿಸಬೇಕು.  ಇದು ಚಿತ್ರ \ref{hdr_init1}ನ ಎಡಭಾಗದಲ್ಲಿ ತೊರಿಸಿದಂತೆ ಕಾಣುತ್ತದೆ. ಈ ಆಜ್ಞೆಯನ್ನು ನಾವು ಬರೆದಿರುವ ಪ್ರೋಗ್ರಾಂನಲ್ಲಿನ "ಒತ್ತಿದಾಗ" ಆಜ್ಞೆಯ ನಂತರ ಜೋಡಿಸಿ,  ಪ್ರೋಗ್ರಾಂ ಅನ್ನು  ಚಲಾಯಿಸಿ ಚೌಕ ಮೂಡಿಬರುವುದನ್ನು ಗಮನಿಸಿ.  ಹಳಯೆ ಚೌಕಗಳೆಲ್ಲ ಅಳಿಸಿ ಒಂದೇ ಚೌಕ ಮೂಡಿಬಂದಿರಬೇಕು. 

ಈಗ ಪ್ರೋಗ್ರಾಂ ಅನ್ನು  ಮತ್ತೊಮ್ಮೆ ಚಲಾಯಿಸಿ ಏನಾಗುತ್ತದೆಂದು ಗಮನಿಸಿ. ಈಗ ಒಂದೇ ಚೌಕ ಮೂಡಿಬಂತಾದರೂ,  ಮೊದಲ್ಲಿದೆಡೆಯಲ್ಲಿ ಬರಲ್ಲಿಲ್ಲ.  ಏಕೆ? ನಾವು ಹಾಕಿದ ಪ್ರೋಗ್ರಾಂ ಶೀರ್ಷಿಕೆ, ವೇದಿಕೆಯನ್ನು ಸ್ವಚ್ಛಗೊಳಿಸಿದೆವಾದರೂ, ಸ್ಪ್ರಯ್ಟ್ ಮೊದಲಿದ್ದೆಡೆಗೆ ಬರಲಿಲ್ಲ. ಇದನ್ನು ತಿದ್ದಲು, ಸ್ಕ್ರಾಚ್ನ "ಚಲನೆ" ಗುಂಪಿನಲ್ಲಿ  \textenglish{x: 0 y: 0} ಗೆ ಹೋಗು ಎಂಬ ಆಜ್ಞೆಯನ್ನು ಉಪಯೋಗಿಸಬೇಕು.  ಚಿತ್ರ \ref{hdr_init1}ನ ನಡುವಲ್ಲಿ ತೊರಿಸಿದಂತೆ  ಪ್ರೋಗ್ರಾಂ ಶೀರ್ಷಿಕೆಯನ್ನು ಬದಲಾಯಿಸಿ ನೋಡಿ. 

ಈಗ ಪ್ರೋಗ್ರಾಂ ಅನ್ನು  ಮತ್ತೊಮ್ಮೆ ಚಲಾಯಿಸಿ ಏನಾಗುತ್ತದೆಂದು ಗಮನಿಸಿ. ಈಗಲೂ ಒಂದೇ ಚೌಕ ಮೂಡಿಬಂತಾದರೂ,  ಮೊದಲ್ಲಿದೆಡೆಯಲ್ಲಿ ಬರಲ್ಲಿಲ್ಲ.  ಏಕೆ? ನಾವು ಹಾಕಿದ ಪ್ರೋಗ್ರಾಂ ಶೀರ್ಷಿಕೆ, ವೇದಿಕೆಯನ್ನು ಸ್ವಚ್ಛಗೊಳಿಸಿ, ಸ್ಪ್ರಯ್ಟ್ ಮೊದಲಿದ್ದೆಡೆಗೆ ತಂದಿಟ್ಟಿತಾದರೂ,  ಸ್ಪ್ರಯ್ಟ್ ಬಲಗಡೆಗೆ ಮುಖಮಾಡಿ ನಿಂತಿರಲಿಲ್ಲ. ಇದನ್ನು ತಿದ್ದಲು, ಸ್ಕ್ರಾಚ್ನ "ಚಲನೆ" ಗುಂಪಿನಲ್ಲಿ  \textenglish{90} ಯ ದಿಕ್ಕಿನಲ್ಲಿಡು ಎಂಬ ಆಜ್ಞೆಯನ್ನು ಉಪಯೋಗಿಸಬೇಕು.  ಚಿತ್ರ \ref{hdr_init1}ನ ಬಲಭಾಗದಲ್ಲಿ ತೊರಿಸಿದಂತೆ  ಪ್ರೋಗ್ರಾಂ ಶೀರ್ಷಿಕೆಯನ್ನು ಬದಲಾಯಿಸಿ, ಪ್ರೋಗ್ರಾಮನ್ನು  ಮತ್ತೊಮ್ಮೆ ಚಲಾಯಿಸಿ ನೋಡಿ. 

ಇಷ್ಟಾದರೂ ಎಲ್ಲಾ ಪ್ರೋಗ್ರಾಮಿನಲ್ಲಿಯೂ ಈ ಶೀರ್ಷಿಕೆಯನ್ನು ಉಪಯೋಗಿಸಲಾಗುವುದಿಲ್ಲ.  ಉದಾಹರಣೆಗೆ ಚಿತ್ರ \ref{hdr_diff}ನಲ್ಲಿನ ಎಡಭಾಗದಲ್ಲಿರುವಂತೆ ಪ್ರೋಗ್ರಾಮ್ಮನ್ನು ಮತ್ತೆ ಮೂಡಿಸಿ ಪ್ರೋಗ್ರಾಂ ಅನ್ನು  ಚಲಾಯಿಸಿ ನೋಡಿ.  ಚೌಕ ಕಾಣಿಸುವುದೇ? ಪ್ರೋಗ್ರಾಮಿನಲ್ಲಿ ಲೇಖನಿಯುಕ್ತ ಎಂಬ ಆಜ್ಞೆಯನ್ನು ಬಳಸದಿದ್ದರೂ ಚೌಕ ಕಾಣಿಸಿತೇಕೆ? ಹಿಂದಿನ ಪ್ರೋಗ್ರಾಂನಲ್ಲಿ ಲೇಖನಿಯುಕ್ತ ಮಾಡಿದನ್ನು ಕಂಪ್ಯೂಟರ್ ನೆನಪಿನಲ್ಲಿಟ್ಟುಕೊಂಡಿರುವುದೇ ಇದಕ್ಕೆ ಕಾರಣ.  ಇದನ್ನು ತಿದ್ದಲು ನಮ್ಮ ಪ್ರೋಗ್ರಾಂ ಶೀರ್ಷಿಕೆಯಲ್ಲಿ ಚಿತ್ರ \ref{hdr_diff}ನಲ್ಲಿನ ಬಲಭಾಗದಲ್ಲಿ ತೋರಿಸಿರುವಂತೆ ಲೇಖನಿಮುಕ್ತ ಎಂಬ ಆಜ್ಞೆಯನ್ನು ಬಳಸಬೇಕು. ಒಮ್ಮೊಮ್ಮೆ ಲೇಖನಿಯ ಬಣ್ಣ, ಗಾತ್ರಗಳನ್ನೂ ಶೀರ್ಷಿಕೆಯಲ್ಲಿ ಸರಿಪಡಿಸ ಬೇಕಾಗಬಹುದು.

\begin{figure}[h]
\begin{multicols}{2}
\begin{center}
\begin{Scratch}[1]
\beginbox{}
\scbox{\cb[w]{100} ಹೆಜ್ಜೆ ಮುಂದೆ ಹೋಗು}{motion}
\turnbox{1}{90}{ಡಿಗ್ರಿಯಷ್ಟು ತಿರುಗು}
\scbox{\cb[w]{100} ಹೆಜ್ಜೆ ಮುಂದೆ ಹೋಗು}{motion}
\turnbox{1}{90}{ಡಿಗ್ರಿಯಷ್ಟು ತಿರುಗು}
\scbox{\cb[w]{100} ಹೆಜ್ಜೆ ಮುಂದೆ ಹೋಗು}{motion}
\turnbox{1}{90}{ಡಿಗ್ರಿಯಷ್ಟು ತಿರುಗು}
\scbox{\cb[w]{100} ಹೆಜ್ಜೆ ಮುಂದೆ ಹೋಗು}{motion}
\end{Scratch}

\begin{Scratch}[1]
\scbox{ಅಳಿಸು}{pen}
\scbox{ಲೇಖನಿಮುಕ್ತ}{pen}
\scbox{\textenglish{x:} \cb[w]{0} \textenglish{y:}\cb[w]{0}ಗೆ ಹೋಗು}{motion}
\scbox{\cb[w]{90} ಯ ದಿಕ್ಕಿನಲ್ಲಿಡು}{motion}
\end{Scratch}
\end{center}
\end{multicols}
\caption{ಸ್ಕ್ರಾಚ್ ಪ್ರೋಗ್ರಾಂ ಹಂತ ಪ್ರಕ್ರಿಯೆ}
\label{hdr_diff}
\end{figure}

\section{ಪ್ರೋಗ್ರಾಂ ಉಳಿಸುವುದು}
\begin{figure}[h]
\begin{center}
\includegraphics[scale=2]{Scratchlogo.pdf}
\end{center}
\caption{ಸ್ಕ್ರಾಚ್ ಲೊಗೊ}
\label{hdr_logo}
\end{figure}

ಕೆಲವು ಪ್ರೋಗ್ರಾಂಗಳನ್ನು ಮುಂದೆಂದಾದರೂ ಬಳಸಲು ಅವನ್ನು ಕಂಪ್ಯೂಟರ್ನಲ್ಲಿಯೆ ಉಳಿಸಬಹುದು. ಉಳಿಸಿದ ಪ್ರೋಗ್ರಾಂಗಳನ್ನು ಸ್ಕ್ರಾಚ್ನಲ್ಲಿ ತೆರೆದು ಪುನಃ ಚಲಾಯಿಸಬಹುದು.   ಪ್ರೋಗ್ರಾಂಗಳನ್ನು  ಉಳಿಸಲು ಸ್ಕ್ರಾಚ್ ಸ್ಕ್ರೀನಿನ ಮೇಲಿನ ಎಡಭಾಗದಲ್ಲಿ ಚಿತ್ರ \ref{hdr_logo}ನಲ್ಲಿ ತೋರಿಸಿರುವಂತೆ ಸ್ಕ್ರಾಚ್ ಲೋಗೋ ಇದೆ. ಈ ಲೋಗೋ ಹತ್ತಿರದ \textenglish{File} ಅನ್ನು ಒತ್ತಿ ಅದರಲ್ಲಿ "ಯೋಜನೆಯನ್ನು ಉಳಿಸು" ಎಂದಿರುವಲ್ಲಿ ಒತ್ತಿದರೆ, ಸಣ್ಣ ಫಾರ್ಂ ಒಂದು ತೆರೆಯುತ್ತದೆ.  ಇದರಲ್ಲಿ ನಮಗೆ ಬೇಕಾದ ಹೆಸರಿನಲ್ಲಿ ನಮ್ಮ ಪ್ರೋಗ್ರಾಮನ್ನು ಉಳಿಸಬಹುದು.  ಇಲ್ಲಿ ಹೆಸರನ್ನು ಕೊಟ್ಟು ಸರಿ ಎಂದು ಒತ್ತಿ ನೋಡಿ. 

ಇದೇ ರೀತಿ ನಾವು ಮಾಡಿರುವ ಪ್ರೋಗ್ರಾಮನ್ನು ತೊರೆದು ಹೊಸ ಪ್ರೋಗ್ರಾಮನ್ನು ಮೂಡಿಸಬಹುದು.  ಪ್ರೋಗ್ರಾಮಿನಲ್ಲಿ ಅಥವಾ ಪ್ರೋಗ್ರಾಂ ಶೀರ್ಷಿಕೆಯಲ್ಲಿ ತಪ್ಪಾದಾಗಲೂ ಹೀಗೆಯೇ  ಪ್ರೋಗ್ರಾಮನ್ನು ತೊರೆದು ಹೊಸ ಪ್ರೋಗ್ರಾಮನ್ನು ಮೂಡಿಸಬಹುದು. \textenglish{File} ಅನ್ನು ಒತ್ತಿ ಅದರಲ್ಲಿ "ಹೊಸ" ಎಂಬಲ್ಲಿ ಒತ್ತಿ, ಈಗ ಉಳಿಸುವುದು "ಬೇಡ" ಎಂದು ಒತ್ತಿದರೆ ಹಳೆಯ ಪ್ರೋಗ್ರಾಂ ಬದಲು ಹೊಸ ಪ್ರೋಗ್ರಾಂಗಾಗಿ ಸ್ವಚ್ಛ ವೇದಿಕೆ ಕಾಣಿಸುತ್ತದೆ. 

\section{ಪ್ರೋಗ್ರಾಂ ಹಂತ ಪ್ರಕ್ರಿಯೆ}
ವಿಭಾಗ \ref{wait} ರಲ್ಲಿ ಕಂಪ್ಯೂಟರ್ ಆಜ್ಞೆಗಳನ್ನು ಅತಿವೇಗವಾಗಿ ಮಾದುವುದರಿಂದ ಪ್ರೋಗ್ರಾಮನ್ನು ವಿಷ್ಲೇಶಿಸಲು ಹಂತಹಂತವಾಗಿ ಕಾಯಿಸುವುದನ್ನು ನೋಡಿದೆವು. ಸ್ಕ್ರಾಚ್ನಲ್ಲಿ ಪ್ರೋಗ್ರಾಮನ್ನು ವಿಷ್ಲೇಶಿಸಲು ಮತ್ತೊಂದು ವಿಧಾನ ಇದೆ. \textenglish{File}  ಪಕ್ಕದಲ್ಲಿ "ಚಿತ್ರಸಂಕಲನ"ಎಂಬಲ್ಲಿ ಒತ್ತಿ, "ಪ್ರತಿಯೊಂದು ಹಂತವನ್ನು ಚಲಿಸು"ಅನ್ನುವುದನ್ನು ಆಯ್ಕೆ ಮಾಡಿ ಪ್ರೋಗ್ರಾಂ  ಹಂತಹಂತವಾಗಿ ಮುಂದುವರೆಯುವುದನ್ನು ಗಮನಿಸಿ.  ಚಿತ್ರಸಂಕಲನದಿಂದ "ಪ್ರತಿಯೊಂದು ಹಂತವನ್ನು ತೋರಿಸು " ಎಂಬಲ್ಲಿ ಒತ್ತಿ, ನಮಗೆ ಸೂಕ್ತವಾದ ವೇಗವನ್ನು ಆಯ್ಕೆ ಮಾಡಬಹುದು.

\begin{figure}[t]
\begin{center}
\begin{Scratch}[1]
\beginbox{}
\scbox{\cb[w]{100} ಹೆಜ್ಜೆ ಮುಂದೆ ಹೋಗು}{motion}
\turnbox{1}{90}{ಡಿಗ್ರಿಯಷ್ಟು ತಿರುಗು}
\scbox{\cb[w]{100} ಹೆಜ್ಜೆ ಮುಂದೆ ಹೋಗು}{motion}
\turnbox{1}{90}{ಡಿಗ್ರಿಯಷ್ಟು ತಿರುಗು}
\scbox{\cb[w]{100} ಹೆಜ್ಜೆ ಮುಂದೆ ಹೋಗು}{motion}
\turnbox{1}{90}{ಡಿಗ್ರಿಯಷ್ಟು ತಿರುಗು}
\scbox{\cb[w]{100} ಹೆಜ್ಜೆ ಮುಂದೆ ಹೋಗು}{motion}
\end{Scratch}
\end{center}
\caption{ಸ್ಕ್ರಾಚ್ ಪ್ರೋಗ್ರಾಂ ಹಂತ ಪ್ರಕ್ರಿಯೆ}
\label{hdr_sstep}
\end{figure}
